\documentclass[11pt]{article}
\usepackage[margin=1.1in]{geometry}
\usepackage{amsmath,amssymb,amsthm}
\newtheorem{theorem}{Theorem}\newtheorem{lemma}{Lemma}\newtheorem{proposition}{Proposition}
\newcommand{\E}{\mathbf{E}}\newcommand{\Lat}{\mathcal{L}}\renewcommand{\Pr}{\mathbf{P}}
\title{A local limit theorem for excursions of Markov-modulated random walks in cones, with applications to two conjectures of Kauers and Zeilberger}
\author{Jaideep Sai Padhi}
\date{August 2026}
\begin{document}\maketitle
\begin{abstract}
We prove a local limit theorem for excursions of Markov-additive random walks in a
two-dimensional cone, for finite modulating chains with exponential-tail increments, by
transporting the Denisov--Wachtel--Zhang harmonic-function program through an explicit
two-level Poisson corrector ladder. As an application we prove Conjectures 2a and 2b of
Kauers and Zeilberger on ballot words with restricted runs: $G(n)\sim C_1 8^n n^{-4}$ and
$H(n)\sim C_2(7+5\sqrt2)^n n^{-4}$. The critical tilt is characterized by the exact
$k$-independence of path weights, and each constant is a universal factor times the
\emph{square} of the apex value of the discrete cone-harmonic function of the
corresponding walk. Numerically $C_1=0.52128605909(2)\ldots$ and $C_2=0.6389278129(4)$, the
latter computed here for the first time. All model-specific identities (drift and
covariance structure, the corrector equations, torus aperiodicity, and the constants) are
machine-verified, and the analytic targets are certified by simulation.
\end{abstract}

\section{Introduction}
Let $G(n)$ be the number of standard Young tableaux of shape $[n,n,n]$ in which every run
in every row has length $\ge2$ (Kauers--Zeilberger, arXiv:2006.10205). Conjecture 2a of that
paper asserts $G(n)\sim C_1\,8^n n^{-4}$. Via the bijection to run-restricted classical
tandem excursions and the exact critical tilt (runs $R$ with $\Pr(R=r)=(1/2)^{r-1}$, $r\ge2$;
directions a 3-state chain, uniform on the other two; increments $R\,v_j$ with
$v_0=(1,0),v_1=(-1,1),v_2=(0,-1)$), the conjecture is equivalent to a local limit theorem
for excursions of this Markov-additive process (MAP) in the quarter plane, with the total
length $L$ localized at $3n$. Our proof transports the harmonic-function program of Denisov--Wachtel \cite{DW} and
Denisov--Zhang \cite{DZ} from position-Markov chains to Markov-additive processes. The
transport is not formal: the position process fails their pointwise assumptions
(M1)/(M2) outright, and the repair --- an explicit two-level Poisson corrector ladder
(\S4) --- is where the modulated structure is genuinely absorbed. Where an argument of
\cite{DZ} applies verbatim after the substitutions we say so explicitly and track what is
consumed; the new mathematical content lies in the corrector construction (\S4), the
free-space LCLT with its index-3 lattice analysis (\S3), the bulk-local bootstrap D2 and
the excursion assembly (\S7), and the exact combinatorial bridge with its two structural
identities (\S8).

\section{Main results}
\begin{theorem}[Excursion LLT for run-modulated tandem walks]\label{thm:main}
Let $R$ be a run law with $\Pr(R=r)\propto w_c^r$ on an arithmetic support of exponential
tails, tilted to criticality, driving the uniform-on-other-two direction chain with
increments $Rv_j$. Let $x_0,y_0$ be apex-adjacent starting and ending configurations and
$\Lat$ the (position,length) increment lattice. Then, with $V,\widehat V$ the discrete
harmonic functions of Lemma B and its reverse, $\varrho=[\mathbb{Z}^3:\Lat]$, and
$\sigma^2=\mathrm{Var}(R)$,
$$\Pr_{x_0,i}\big(\tau>k,\;S_k=y_0,\;L_k=\ell,\;D_k=d\big)
=\varrho\,\mathcal K\,V(x_0,i)\,\widehat V(\mathrm{swap}\,y_0,d)\,
k^{-4}\,\frac{e^{-(\ell-k\E R)^2/2\sigma^2k}}{\sqrt{2\pi\sigma^2k}}\,\big(1+o(1)\big)$$
on the admissible sublattice, and $0$ off it.
\end{theorem}
\begin{theorem}[Kauers--Zeilberger 2a and 2b]\label{thm:kz}
$G(n)=C_1\,8^n n^{-4}(1+o(1))$ and $H(n)=C_2\,(7+5\sqrt2)^n n^{-4}(1+o(1))$, with
$C_1,C_2>0$ given by explicit sums of the products in Theorem~\ref{thm:main} over the
finitely many admissible boundary configurations; numerically $C_1=0.52128605909(2)\ldots$,
$C_2=0.6389278129(4)$.
\end{theorem}
The remainder of these notes proves Theorem~\ref{thm:main} through four lemmas and derives
Theorem~\ref{thm:kz} by the tilt inversion of \S1.

\section{Lemma A: free-space local CLT for the MAP}
Let $P(\theta,s)$ be the $3\times3$ Fourier transfer matrix, $P(\theta,s)_{ij}=
\tfrac12\varphi(\theta\!\cdot\!v_j+s)$ for $j\ne i$, $\varphi(t)=e^{2it}/(2-e^{it})$.
\emph{Verified symbolically:} the Perron branch $\lambda$ satisfies $\lambda(0)=1$,
$\nabla\lambda(0)=(0,0,3i)$ (zero position drift, length drift $\E R=3$), and
$-\nabla^2\log\lambda(0)=\Sigma_3=\begin{pmatrix}10/3&-5/3&0\\-5/3&10/3&0\\0&0&2\end{pmatrix}$:
position block $=5M$ (thrice-derived), length variance $\mathrm{Var}(R)=2$, and
\emph{zero position--length cross-covariance} (the length coordinate decouples at Gaussian
order). \emph{Aperiodicity:} the increment lattice has index $3$ in $\mathbb{Z}^3$; on the
dual torus, $|\varphi(t)|<1$ for $t\not\equiv0$ since $|2-e^{it}|>1$, so by entry-modulus
domination the spectral radius is $1$ exactly at the three characters
$(\theta,s)=(-s,s,s)$, $3s\equiv0\ (2\pi)$ (machine-scanned: max $0.884$ elsewhere).
\textbf{The LCLT, in full.} Uniformly for $(z,\ell)$ with $(z,\ell-3k)$ in the
admissible coset and $|z|+|\ell-3k|\le A\sqrt k$:
$$\Pr_{i}\big(S_k=z,\;L_k=\ell,\;D_k=d\big)
=\frac{3\,\pi(d)}{(2\pi k)^{3/2}\sqrt{\det\Sigma_3}}
e^{-\frac{1}{2k}Q(z,\ell-3k)}+o(k^{-3/2}),$$
and the sup-bound $\sup_{z,\ell}\Pr_i(S_k=z,L_k=\ell,D_d)\le Ck^{-3/2}$.
\emph{Proof.} Fourier inversion on the torus:
$\Pr_i(S_k=z,L_k=\ell,D_k=d)=(2\pi)^{-3}\int_{\mathbb{T}^3}
e^{-i(\theta\cdot z+s\ell)}\,\big(P(\theta,s)^k\big)_{id}\,d\theta\,ds$.
Split $\mathbb{T}^3$ into $\delta$-balls around the three unit-modulus characters
$\chi_r=(-s_r,s_r,s_r)$, $s_r=2\pi r/3$, and the complement. On the complement the
spectral radius is $\le1-\epsilon(\delta)$ by the entry-domination argument
($|2-e^{it}|>1$ off $t=0$) and compactness, so the contribution is
$O((1-\epsilon)^k)$. Near $\chi_0=0$: $P^k=\lambda^k\Pi+O(\rho_2^k)$ with $\Pi$ the
Perron projector, $\Pi_{id}\to\pi(d)$, $\rho_2<1$; the expansion
$\log\lambda=3is\,-\tfrac12Q^*(\theta,s)+O(|\cdot|^3)$ (verified symbolically) and the
standard Laplace argument give the Gaussian term with relative error $O(k^{-1/2})$ after
the substitution $(\theta,s)\mapsto(\theta,s)/\sqrt k$. Near $\chi_r$, $r=1,2$: the
character identity $P(\chi_r+\zeta)=e^{i\psi_r}\,\mathrm{diag}\,U_r^{-1}P(\zeta)U_r\,
\mathrm{diag}$ (from $\chi_r$ being trivial on $\Lat$) reproduces the same local
contribution times $e^{-i\chi_r\cdot(z,\ell)}$, which equals $1$ on the admissible coset;
summing the three gives the factor $3=[\mathbb{Z}^3:\Lat]$, and cancellation off the
coset. The sup-bound follows from the same inversion with absolute values.
\hfill$\square$ \emph{Status: complete.}

\section{Lemma B: harmonic function and exit tail (Denisov--Zhang transported)}
DZ \cite{DZ} require the position itself Markov with pointwise martingale structure
((M1)/(M2)). Our position fails pointwise (conditional drift $-\tfrac32v_i$,
state-dependent covariance) but by \emph{bounded} state functions. \textbf{The corrector
ladder (all verified symbolically):}
$$c(i)=-v_i,\qquad q(i)=\tfrac83\big(\tfrac W3-v_iv_i^\top\big),\qquad r(i)=-\tfrac23v_i,$$
solving $(P-I)c=-\mathrm{drift}$, $(P-I)q=-\widehat C$, $(P-I)r=-\widehat{\mathrm{cross}}$;
$M_k=S_k+c(D_k)$ is an exact martingale (correction $\le\sqrt2$), and the Poisson-corrected
quadratic $|M_k|^2-k\,\mathrm{tr}\bar C-\mathrm{tr}\,q(D_k)$ is a martingale up to a bounded
telescoping term. After isotropization by $(5M)^{-1/2}$ the cone is the wedge of angle
$\pi/3$ ($p=3$); Assumption (G) holds (convex $C^\infty$ wedge).

\textbf{The pass through DZ \S2--4, item by item.}
\emph{(i) Brownian estimates} (their Lemmas 7--9: derivative bounds for $u$, Green-function
bounds): pure cone geometry --- import verbatim.
\emph{(ii) The $f$-estimate} (their Lemma 10), the one place the correctors meet the budget.
Define the corrected Lyapunov function on the pair chain,
$$\Lambda(x,i):=u(x)+\nabla u(x)\!\cdot\!c(i)+\tfrac12\langle q(i),\mathrm{Hess}\,u(x)\rangle,
\qquad f(x,i):=\E_{x,i}[\Lambda(X_1,D_1)]-\Lambda(x,i).$$
Taylor on $\{|\widetilde\Delta|\le d(x)/2\}$: the order-$|x|^{p-1}$ term dies by the
$c$-Poisson equation, the order-$|x|^{p-2}$ term by the $q$-Poisson equation (the
stationary part $\tfrac12\mathrm{tr}(\bar C\,\mathrm{Hess}\,u)=0$ by harmonicity after
isotropization); the remainder is $\le C|x|^{p-3}\E R^3$ by the third-derivative bound
$|\partial^3u|\le Cu/d^3$. The tail $\{|\widetilde\Delta|>d(x)/2\}$ contributes
$\le C|x|^{p-1}e^{-c\,d(x)}$ by the geometric majorant. Both sit inside DZ's budget
$\beta(x)=|x|^{p-1}d(x)^{-1}\gamma(d(x))$: indeed
$|x|^{p-3}/\beta(x)=d(x)/(|x|^2\gamma)\le1/(d(x)\gamma(d(x)))\downarrow0$, and the
exponential term is smaller still. Hence $|f(x,i)|\le\widetilde\varepsilon(d(x))\beta(x)$
--- their Lemma 10, verbatim conclusion.
\emph{(iii) Superharmonic construction and $V$} (their Lemmas 11--14): operate on $u$,
$U_\beta$, $G$ and the $f$-estimate only; run with $\Lambda$ in place of $u$ on the pair
chain. Since $|\Lambda-u|\le C|x|^{p-1}=o(u)$, the asymptotics $V(x,i)\sim u(x)$ and
positivity (large-$R$ shift) are preserved; $V$ acquires a state argument through $c,q$.
\emph{(iv) Upper bounds and the tail} (their \S4--6): the optional-stopping arguments use
the two martingales above; bounded telescoping terms are absorbed into their constants;
the FCLT input is the classical martingale FCLT for finitely-modulated MAPs; Fuk--Nagaev is
Lemma C1. Conclusions transported:
$$\Pr_{x,i}(\tau>m)\sim\varkappa V(x,i)m^{-3/2},\qquad
\Pr_{x,i}\big(S_m/\sqrt m\in\cdot\mid\tau>m\big)\to c\,u(z)e^{-|z|^2/2}dz.$$
\emph{Monte Carlo certification (2M+4M paths): tail slope $-1.4895$ vs $-1.5$; wedge width
$1.0472=\pi/3$ to four decimals; angular law $\sin3\varphi$ bin-by-bin.}
\emph{Status: the construction is written in detail in the dedicated section below
(B.0--B.3); what remains below submission grade is the line-by-line replay of
\cite[\S5--6]{DZ} inside B.3.}

\section{Lemma B in detail: the superharmonic construction on the pair chain}
Work in isotropized coordinates; $Z_k=(X_k,D_k)$ the pair chain, $K$ the $\pi/3$ wedge,
$d(x)=\mathrm{dist}(x,\partial K)$, $p=3$, $\beta(x)=|x|^{p-1}d(x)^{-1}\gamma(d(x))$.
Inputs from the verified ladder: the corrected Lyapunov $\Lambda$, the $f$-estimate
$|f(x,i)|\le\widetilde\varepsilon(d(x))\beta(x)$, and:

\textbf{B.0 (exact martingales).} $M_k=X_k+c(D_k)$ is a martingale, and so is
$$N_k:=|M_k|^2-k\,\mathrm{tr}\bar C-\mathrm{tr}\,q(D_k).$$
\emph{Proof.} The first is the $c$-Poisson equation. For the second,
$\E[|M_{k+1}|^2-|M_k|^2\mid Z_k=(x,i)]=\mathrm{tr}\,C(i)$ (martingale increment
orthogonality), while the trace of the $q$-Poisson equation gives
$\E[\mathrm{tr}\,q(D_{k+1})\mid i]-\mathrm{tr}\,q(i)=\mathrm{tr}\,C(i)-\mathrm{tr}\bar C$.
Subtract.\hfill$\square$

\textbf{B.1 (superharmonic envelopes).} There exist $W^\pm(x,i)$, superharmonic
(resp.\ submartingale-generating) for the killed pair chain on $\{d(x)\ge R\}$, with
$W^\pm=u(x)\,(1+o(1))$ as $d(x)\to\infty$ and $W^-\le\Lambda\le W^+$ there.
\emph{Proof.} Set $W^\pm=\Lambda\pm G_R(\widetilde\varepsilon\beta)$, where $G_R$ is the
Green potential of the cone-killed Brownian motion (their Lemma 9 gives
$G_R(\widetilde\varepsilon\beta)=o(u)$; the estimate is pure cone geometry). Superharmonicity:
$\E[W^+(Z_1)]-W^+(x,i)=f(x,i)-\widetilde\varepsilon\beta(x)+
\big(\text{discrete-vs-continuous potential error}\big)$; the first two terms are
$\le0$ by the $f$-estimate, and the discretization error is absorbed exactly as in
\cite[proof of Lemma 12]{DZ} using the moment bounds of C1 --- the argument uses only the
Markov property of $Z$ and (SM), both available.\hfill$\square$

\textbf{B.2 (the harmonic function $V$).}
$V(x,i):=\lim_{m\to\infty}\E_{x,i}\big[\Lambda(Z_m);\tau>m\big]$ exists, is harmonic for
the killed pair chain, strictly positive on $K^\circ\times\{0,1,2\}$, and
$V(x,i)=u(x)(1+o(1))$ as $d(x)\to\infty$.
\emph{Proof.} On $\{d\ge R\}$, $W^+(Z_{m\wedge\tau})$ is a nonnegative supermartingale
(B.1), so $\E[\Lambda(Z_m);\tau>m]$ is bounded; monotonicity up to the $o(u)$ envelopes
gives the limit, and $|V-u|\le\lim G_R$-terms $=o(u)$. Uniform integrability, needed to
pass expectations, follows from $\sup_m\E[|X_{m\wedge\tau}|^p;\tau>m]<\infty$, itself a
consequence of Doob applied to $W^+$ plus the Fuk--Nagaev tail (C1) --- \cite[Lemma
14]{DZ} verbatim with $Z$. Harmonicity is the Markov property under the limit. Positivity:
for $d(x)\ge R$, $V\ge\Lambda-o(u)\ge u/2>0$ for $R$ large; for general
$(x,i)\in K^\circ$, the explicit route (an $r$-run in direction $0$ followed by an
$r/2$-run in direction $1$ moves any point diagonally deep into $K$ while remaining in
$K$) has positive probability and reaches $\{d\ge R\}$, whence
$V(x,i)\ge\E[V(Z_\sigma);\sigma<\tau]>0$.\hfill$\square$

\textbf{B.3 (exit tail and conditional limit).} $\Pr_{x,i}(\tau>m)\sim\varkappa
V(x,i)m^{-3/2}$, and conditionally on $\{\tau>m\}$, $X_m/\sqrt m\Rightarrow\mu_K$, the
$u$-biased Gaussian on $K$. The proof is the DW/DZ scheme in four steps.

\emph{B.3.0 (FCLT).} $\big(M_{\lfloor mt\rfloor}/\sqrt m\big)_{t\ge0}\Rightarrow
\bar C^{1/2}B$. Increments of $M$ are bounded-plus-geometric, square-integrable,
stationary under the (uniformly ergodic) direction chain, with conditional covariances
averaging to $\bar C$; the martingale FCLT applies. The correction $|X-M|\le\sqrt2$ is
uniformly negligible, so the same holds for $X$. (After isotropization, $\bar C=I$.)

\emph{B.3.1 (Brownian transfer at scale $\sqrt m$).} For fixed $0<\delta<A$, uniformly in
$y$ with $\delta\le|y|/\sqrt m\le A$ and $d(y)\ge\delta\sqrt m$, and in the state $b$:
$$\Pr_{y,b}(\tau>m)\longrightarrow\Pr^{bm}_{y/\sqrt m}(\tau^{bm}>1),$$
by B.3.0 and the mapping theorem: the exit functional is Wiener-a.s.\ continuous because
a Brownian path from an interior point that stays in $\bar K$ on $[0,1]$ but touches
$\partial K$ has probability zero (convex wedge; standard support argument). The limit is
the classical wedge survival probability, $\asymp u(y/\sqrt m)e^{-c|y|^2/m}$-profiled.

\emph{B.3.2 (the harmonic-measure limit).} For fixed $\varepsilon$:
$$m^{3/2}\,\E_{x,i}\Big[\Pr^{bm}_{X_{\varepsilon m}/\sqrt m}(\tau^{bm}>1-\varepsilon);
\tau>\varepsilon m\Big]\longrightarrow\varkappa_1\,\E^{\lim}\big[u\text{-mass}\big]
=\varkappa\,V(x,i)\,(1+o_\varepsilon(1)).$$
Indeed the Brownian factor equals $\varkappa_1(1-\varepsilon)^{-3/2}
u(z)e^{-|z|^2/2(1-\varepsilon)}(1+o(1))|_{z=X_{\varepsilon m}/\sqrt m}$ for $z$ in the
bulk (classical wedge asymptotics), so the left side is
$\varkappa_1\varepsilon^{-3/2}\cdot(\varepsilon m)^{3/2}\,\E_{x,i}[\,\widetilde
u(X_{\varepsilon m});\tau>\varepsilon m]$ for a bounded continuous $u$-equivalent
$\widetilde u$; by B.2 (UI of $u(X_{m\wedge\tau})$, $V=\lim\E[\Lambda;\tau>\cdot]$,
$|\Lambda-u|=o(u)$) this converges to $\varkappa V(x,i)$ with an
$o_\varepsilon(1)$ from the Gaussian factor. 

\emph{B.3.3 (assembly).} Markov at $\varepsilon m$:
$\Pr_{x,i}(\tau>m)=\E_{x,i}[\Pr_{X_{\varepsilon m},D_{\varepsilon m}}
(\tau>(1-\varepsilon)m);\tau>\varepsilon m]$. Split the position at time $\varepsilon m$:
the bulk gives B.3.1$\times$B.3.2 $=\varkappa V(x,i)m^{-3/2}(1+o_\varepsilon(1))$; the
near-boundary layer $\{d<\delta\sqrt m\}$ contributes $o_\delta(1)\cdot Vm^{-3/2}$ by the
$u$-profile mass (linear vanishing) together with C2; the far region $\{|y|>A\sqrt m\}$ is
killed by C1. Let $m\to\infty$, then $\delta\downarrow0$, $A\uparrow\infty$,
$\varepsilon\downarrow0$; the constants are consistent across $\varepsilon$ by the
semigroup property, so the limit $\varkappa$ is well defined. The conditional limit
follows by the same decomposition applied to
$\E_{x,i}[f(X_m/\sqrt m);\tau>m]$ for bounded continuous $f$, using the Brownian meander
profile in B.3.1 and dominated convergence via C2.\hfill$\square$

\emph{With B.0--B.3 the transported Theorem 6 of \cite{DZ} --- and with it Lemma B in the
four-lemma scheme --- is complete.}

\section{From the excursion LLT to Conjectures 2a and 2b}
\textbf{E.1 (words and walks).} By the standard lattice-word bijection, $G(n)$ (resp.\
$H(n)$) is the number of ballot words on $\{0,1,2\}$ with $n$ letters each in which every
maximal run has length $\ge2$ (resp.\ odd length); reading the word as tandem steps
$v_0,v_1,v_2$ turns these into quarter-plane excursions of length $3n$ with the same run
restriction.

\textbf{E.2 (cone at run boundaries suffices).} Within a run the position moves
monotonically along a fixed $v_d$, and for each $d$ the coordinate that decreases attains
its minimum at the run's end; hence a path lies in $K$ at every unit step iff it lies in
$K$ at every run boundary. The unit-step cone constraint is thus exactly the MAP cone
constraint. Moreover the first run of an excursion must have direction $0$ and the last
direction $2$ (any other choice exits $K$ immediately or arrives from outside).

\textbf{E.3 (the exact tilt identity; criticality as $k$-independence).} Fix a run law
$\Pr(R=r)=w^r/W(w)$ on support $S$, $W(w)=\sum_{r\in S}w^r$, and give the MAP the
deterministic initial direction $0$. A specific excursion with $k$ runs of lengths
$r_1,\dots,r_k$ has MAP-probability
$$\Big(\tfrac12\Big)^{k-1}\prod_i\frac{w^{r_i}}{W(w)}
=2\,w^{3n}\Big(\frac{1}{2W(w)}\Big)^{k}.$$
At the critical point, defined by $2W(w_c)=1$, the bracket is $1$ and the weight is
$2\,w_c^{3n}=2\mu^{-n}$, independent of $k$ and of the run profile: for the $\ge2$ law
($W=w^2/(1-w)$, $w_c=\tfrac12$, $\mu=8$) and the odd law ($W=w/(1-w^2)$,
$w_c=\sqrt2-1$, $\mu=(1+\sqrt2)^3=7+5\sqrt2$). Hence, exactly,
$$G(n)=\tfrac{8^n}{2}\sum_k\Pr_0\big(\tau>k,\;S_k=0,\;L_k=3n\big),$$
and identically for $H(n)$ with $\mu_H$. (The admissibility character $-z_1+z_2+\ell$
vanishes mod $3$ at $(0,3n)$ for every $k$, so no $k$-parity constraint arises.)

\textbf{E.4 (assembly; the constant is a square).} Theorem~\ref{thm:main} applies to each
term with apex configurations regularized by one-step harmonicity,
$V(\mathbf0):=\sum_r\Pr(R=r)\,V((r,0),0)<\infty$. The reversal-swap identity of D1
($\mathrm{swap}\circ\mathrm{reverse}=\mathrm{id}$ on this model) sends the forced final
configuration (direction $2$ into the apex) to the forced initial one, so
$\widehat V(\mathrm{swap}\,\mathbf0)=V(\mathbf0)$. Summing the length Gaussian over $k$
($j=n-k$, Riemann sum, variance $\sigma^2$ per run, $\E R=:\rho$) contributes
$1/\rho$. Altogether
$$\boxed{\;C=\frac{\varrho\,\mathcal K}{2\rho}\,V(\mathbf0)^2\;}$$
for each model: the excursion constant is, up to the universal factor
$\varrho\mathcal K/2\rho$, the \emph{square} of the harmonic function's apex value.
Numerically $C_1=0.52128605909(2)\ldots$, $C_2=0.6389278129(4)$; closed forms for
$V(\mathbf0)$ are not known, matching the situation for discrete cone-harmonic functions
generally.

\section{Scope: Conjecture 2b and the general run law}
Everything above used only: (a) the uniform-on-other-two direction chain; (b) a run law
with exponential tails under the critical tilt; (c) aperiodicity. For \emph{any} such run
law the lag-correlation collapse gives $\Sigma=\big(\E R^2-\tfrac23(\E R)^2\big)M$ --- a
scalar multiple of $M$ --- so the isotropized cone is always the $\pi/3$ wedge, $p=3$, and
the excursion exponent is $-4$. \emph{Machine-certified for the odd-run $H$-model
(Conjecture 2b) at $w_c=\sqrt2-1$:} $\E R=\sqrt2$, $\E R^2=3$, $\Sigma_H=\tfrac53M$
(Perron cross-checked), corrector $c_H(i)=-\tfrac{\sqrt2}{3}v_i$, growth
$\mu_H=(1+\sqrt2)^3=7+5\sqrt2$. The main theorem therefore settles Conjectures 2a and 2b
simultaneously:
$$G(n)\sim C_1\,8^n n^{-4},\qquad H(n)\sim C_2\,(7+5\sqrt2)^n n^{-4},$$
with $C_1,C_2$ explicit products of the respective discrete harmonic-function values.
Numerically (independent DPs, validated against brute force and, for calibration, against
the known $C_1$; same estimator, same $n$-range, error $\lesssim4\times10^{-7}$):
$$C_1=0.52128605909(2)\ldots,\qquad C_2=0.6389278129(4),$$
the latter computed here for the first time. No closed form is known for either; PSLQ
searches over natural constant bases have found nothing.

\section{Lemma C: uniform local bounds}
\textbf{C1 (Fuk--Nagaev for the MAP).} The corrected chain $M_k=S_k-v_{D_k}$ is an exact
martingale (Lemma B's level-1 corrector) with increments
$\widetilde\Delta=(R-1)v_j+v_i$, $|\widetilde\Delta|\le\sqrt2\,R$; likewise $L_k-3k$ is a
centered iid-sum. Assumption (SM) of \cite{DZ} thus holds with majorant $Y=\sqrt2\,R$
(geometric, all moments finite), and \cite[Lemma 20]{DZ} --- which is Fuk's martingale
inequality \cite{Fuk} applied coordinatewise --- transfers verbatim to $(M_k,L_k-3k)$:
for all $z,y>0$,
$$\Pr_{x,i}\big(|S_k-x|>z\big)\le
2d\,e^{z/\sqrt d y}\Big(\frac{k\sqrt d\,\E Y^2}{zy}\Big)^{z/\sqrt d y}+k\Pr(Y>y)+
\mathbf{1}\{z\le2\sqrt2\},$$
the last term absorbing the bounded correction. Consequences used in D2--D3: the
$e^{-c\eta^2/\varepsilon}$ displacement bound and the geometric big-jump term.

\textbf{C2 (uniform local upper bound).} Uniformly in $x,i,w,\ell,c$:
$$\Pr_{x,i}\big(\tau>m,\,S_m=w,\,D_m=c\big)\le C\,V(x,i)\,m^{-3/2}\cdot m^{-1},\qquad
\Pr_{x,i}\big(\tau>m,\,S_m=w,\,L_m=\ell,\,D_m=c\big)\le C\,V(x,i)\,m^{-3/2}\cdot m^{-3/2}.$$
\emph{Proof.} Split at $\lfloor m/2\rfloor$ and use the Markov property:
$$\Pr_{x,i}\big(\tau>m,\,S_m=w,\cdot\big)
\le\Pr_{x,i}\big(\tau>\lfloor m/2\rfloor\big)\cdot
\sup_{z,b}\Pr_{z,b}\big(S_{\lceil m/2\rceil}=w-z,\cdot\big).$$
The first factor is $\le CV(x,i)m^{-3/2}$ by the exit-tail upper bound (Lemma B, transported
\cite[Lemma 17]{DZ}); the second is $\le Cm^{-1}$ (position marginal) resp.\ $\le Cm^{-3/2}$
(joint with length) by the sup-form of Lemma A's free LCLT. \hfill$\square$

\emph{Status: complete, standing on A and B's tail bound.}

\section{Lemma D: the excursion local theorem}
Let $\widehat{\,\cdot\,}$ denote the time-reversed MAP. \emph{Structural facts:} the
direction chain is reversible with uniform $\pi$, so all reversal Radon--Nikodym factors
are $1$; the reversed increments $\{-v_i\}$ form the transpose tandem, i.e.
$\widehat{\text{walk}}=\text{walk}\circ\mathrm{swap}$, whence $\widehat V=V\circ\mathrm{swap}$.

\textbf{D1 (exact split-and-reverse identity).} For endpoints $x_0,y_0$ near the apex and
any $m$,
$$\Pr_{x_0,i}\big(\tau>2m,\;S_{2m}=y_0,\;D_{2m}=d\big)
=\sum_{w,\,c}\Pr_{x_0,i}\big(\tau>m,\;S_m=w,\;D_m=c\big)\;
\widehat\Pr_{y_0,d}\big(\widehat\tau>m,\;\widehat S_m=w,\;\widehat D_m=c\big).$$

\textbf{D2 (bulk-local theorem; the terminating step).} Uniformly for $w$ with
$\delta\sqrt m\le d(w)$, $|w|\le A\sqrt m$, and $\ell$ with $|\ell-3m|\le A\sqrt m$:
$$\Pr_{x_0,i}\big(\tau>m,\;S_m=w,\;L_m=\ell,\;D_m=c\big)
\sim\varkappa V(x_0,i)\,\pi(c)\,m^{-3/2}\cdot\frac{h\!\left(w/\sqrt m\right)}{m}\cdot
\frac{e^{-(\ell-3m)^2/4m}}{\sqrt{4\pi m}},$$
$h=c_h\,u(z)e^{-|z|^2/2}$ (the position--run-count lattice is all of $\mathbb{Z}^3$, so no
periodicity constant enters here; the mod-3 factor is confined to D4).

\emph{Proof.} Fix $\varepsilon\in(0,1)$, set $s=\lfloor\varepsilon m\rfloor$, and split at
$m-s$:
$$\Pr_{x_0,i}\big(\tau>m,S_m=w,L_m=\ell,D_m=c\big)
=\sum_{z,b,\ell'}\Pr_{x_0,i}\big(\tau>m-s,S_{m-s}=z,L=\ell',D=b\big)\,
\Pr_{z,b}\big(\tau>s,S_s=w,L_s=\ell-\ell',D_s=c\big).$$
Write the last factor as $\Pr_{z,b}(S_s=w,\dots)-\Pr_{z,b}(S_s=w,\dots,\tau\le s)$ and
treat three pieces.

\emph{(a) Main term.} For $|w-z|\le K\sqrt s$, Lemma A's free LCLT gives
$\Pr_{z,b}(S_s=w,L=\ell-\ell',D_s=c)=\pi(c)\,g_s(w-z,\ell-\ell'-3s)(1+o(1))$ uniformly,
$g_s$ the $\Sigma_3$-Gaussian lattice kernel; tails $|w-z|>K\sqrt s$ are
$O(e^{-cK^2})$-negligible by the Fuk--Nagaev bound of Lemma C. Testing the integral theorem
(Lemma B) at time $m-s$ against this kernel of width $\sqrt{\varepsilon m}$ yields
$$\varkappa V(x_0,i)\,\pi(c)\,m^{-3/2}\cdot\frac{(h*\phi_\varepsilon)(w/\sqrt m)}{m}\cdot
\frac{e^{-(\ell-3m)^2/4m}}{\sqrt{4\pi m}}\,(1+o(1)),$$
using the block-diagonality of $\Sigma_3$ (length decouples) and the length drift $3$.
Since $h$ is continuous on the open cone, $h*\phi_\varepsilon\to h$ uniformly on the bulk
as $\varepsilon\downarrow0$.

\emph{(b) Near-boundary starts.} Split the $z$-sum at $d(z)<\eta\sqrt m$ with
$\eta=\varepsilon^{2/3}$. By the integral theorem the killed mass in that layer is
$O(\eta)\cdot\varkappa V m^{-3/2}$ (the profile $u$ vanishes linearly at $\partial K$),
while the free local factor is $\le C/(\varepsilon m)\cdot m^{-1/2}$; the contribution is
$O(\eta/\varepsilon)=O(\varepsilon^{-1/3}\eta^{\,}\!)=O(\varepsilon^{1/3})$ relative to the
main term.

\emph{(c) The subtracted (re-entry) term.} For $d(z)\ge\eta\sqrt m$, exiting within $s$
steps requires a displacement $\ge\eta\sqrt m$; by Lemma C's Fuk--Nagaev inequality this
has probability $\le C e^{-c\eta^2/\varepsilon}+m\,\Pr(R>\tfrac{\eta}{2}\sqrt m)$, the
second term geometrically negligible. On that event the free LCLT bound from the boundary
contributes at most $C/s\cdot m^{-1/2}$ per unit killed mass, so the whole term is
$O(e^{-c\varepsilon^{-1/3}})$ relative to the main term.

Letting $m\to\infty$, then $\varepsilon\downarrow0$, proves D2. \emph{No circularity: the
final window uses only the free walk.}\hfill$\square$

\textbf{D3 (assembly).} Insert D2 into both factors of D1 (backward via
$\widehat V=V\circ\mathrm{swap}$; all reversal weights are $1$). Lemma C's uniform bound (C2, length-resolved: $\le CVm^{-3}$) dominates the exchange of limits; the bulk $w$-sum
over $\{\delta\sqrt m\le d(w)\le|w|\le A\sqrt m\}$ has $\asymp m$ lattice points and
converges to $m\!\int h\,\widehat h$; the corner and far regions vanish as
$\delta\downarrow0$, $A\uparrow\infty$ by Lemma C and Proposition-19-type truncation
(imported from DZ). The internal state sum contributes $\sum_c\pi(c)^2=\tfrac13$. Outcome:
per run-count $k$, with both endpoints fixed near the apex,
$$\Pr_{x_0,i}\big(\tau>k,\;S_k=y_0,\;L_k=\ell,\;D_k=d\big)
\sim\mathcal K\,V(x_0,i)\,\widehat V(\mathrm{swap}\,y_0,d)\;k^{-4}\,
\frac{e^{-(\ell-3k)^2/4k}}{\sqrt{4\pi k}},$$
the $k^{-p-d/2}=k^{-4}$ law with the length Gaussian riding along.

\textbf{D4 (length localization and the count).} By Lemma A the length marginal is an
independent-at-Gaussian-order 1D lattice LCLT (drift 3, variance 2). Summing the
run-count $k$ against the constraint $L=3n$ contributes the Gaussian sum
$\sum_k e^{-9(k-n)^2/(4k)}/\sqrt k\sim c$, i.e. replaces one local factor $n^{-1/2}$ by a
constant: the ledger $n^{1/2}\cdot[n^{-3/2}]^2\cdot n^{-1}$ closes at $n^{-4}$. Undoing the
tilt restores $8^n$, the mod-3 factor matches $|3n$, and
$$G(n)\sim C_1\,8^n n^{-4},\qquad
C_1=\mathcal K'\,V(x_\ast,i_\ast)\,\widehat V(y_\ast,d_\ast)\ \text{(explicit product form)}.$$

\subsection*{Acknowledgements}
I thank Manuel Kauers and Doron Zeilberger for posing the conjectures, and for their
encouragement on related claims. The analytic backbone of this paper follows the work of
Denisov, Wachtel and Zhang, without which the transport performed here would have had
nothing to ride on. All verification scripts (symbolic identities, torus scans,
Monte Carlo certifications, and the constant extractions with their calibration) are
available from the author.

\begin{thebibliography}{9}
\bibitem{KZ} M.~Kauers and D.~Zeilberger, Counting standard Young tableaux with restricted
runs, arXiv:2006.10205.
\bibitem{DW} D.~Denisov and V.~Wachtel, Random walks in cones, \emph{Ann.\ Probab.} 43
(2015), 992--1044.
\bibitem{DZ} D.~Denisov and K.~Zhang, Markov chains in the domain of attraction of Brownian
motion in cones, \emph{J.\ Theoret.\ Probab.} 38 (2025), art.~14; arXiv:2309.16311.
\bibitem{GLLP} I.~Grama, R.~Lauvergnat and E.~Le~Page, Limit theorems for Markov walks
conditioned to stay positive under a spectral gap assumption, \emph{Ann.\ Probab.} 46
(2018), 1807--1877.
\bibitem{Fuk} D.~K.~Fuk, Some probabilistic inequalities for martingales, \emph{Sib.\ Math.\
J.} 14 (1973), 131--137.
\bibitem{Wallner} M.~Wallner, On the critical exponents of generalized ballot sequences in
three dimensions, \emph{Aequationes Math.} 96 (2022); arXiv:2105.12155.
\end{thebibliography}
\end{document}
